\documentclass[11pt,a4paper]{article}
\usepackage[margin=1in]{geometry}
\usepackage{hyperref}
\usepackage{amsmath}
\usepackage[T1]{fontenc}

\hypersetup{
    colorlinks=true,
    linkcolor=blue,
    urlcolor=blue,
    citecolor=blue
}

\setlength{\parindent}{0pt}
\setlength{\parskip}{6pt}

\title{\vspace{-2em}\textbf{Saccade Sequence Parsing:\\Energy-Minimizing Substitution as a Unified\\Mechanism for Visual Scene Structure}\\[0.5em]\large A Proposal}
\author{Generated by Claude Opus 4.6, prompted by Rob Freeman}
\date{March 13, 2026}

\begin{document}
\maketitle
\thispagestyle{empty}

\section{Introduction and Thesis}

We have built a working parser that discovers hierarchical structure in natural language text from raw sequential data, using no grammar rules and no labelled training data. The mechanism is energy-minimizing substitution: for each possible decomposition of a text span into sub-spans, the system finds which words can substitute for each other in context, scores the decomposition by the richness and coherence of these substitution classes, and selects the minimum-energy (best-supported) hierarchical structure. Both the ``objects'' of the representation (substitution classes) and the ``distances'' between them (the energy landscape) emerge from the process itself, rather than being pre-specified.

This paper proposes that the same mechanism applies to visual perception. The key observation is that human vision is not instantaneous---it is \textit{sequential}. We sample visual scenes through a series of saccadic eye movements, producing an ordered sequence of fixation patches. We propose treating these fixation sequences exactly as the parser treats word sequences: building a corpus of fixation subsequences, computing substitution classes (which fixation patches can replace each other in context), and finding the minimum-energy hierarchical decomposition. The resulting structure---visualized as nested containment circles with sequence links threading through---would directly represent the part-whole hierarchy of visual objects.

This proposal connects to Goertzel's recent framework in which evidence functions as energy in logical systems. The parser embodies this principle: the energy of a structural hypothesis is determined entirely by corpus evidence, and evidence conservation prevents the system from ``hallucinating'' structure not grounded in observed data. Extending this to vision means that visual object structure would be grounded in fixation evidence with the same anti-hallucination guarantees.

\section{The Hamiltonian Dynamic Parser: Brief Review}

The parser operates on a text corpus and analyzes input phrases by discovering their optimal hierarchical decomposition. The energy function for each candidate decomposition is:

\[
E = -\log(n_{\text{subs}}) + \log(\text{consensus}) - \log(n_{\text{full\_length}})
\]

where $n_{\text{subs}}$ is the number of valid substitutes found, consensus measures how tightly the substitution class constrains its context (higher consensus = more constrained = higher energy), and $n_{\text{full\_length}}$ counts substitutes matching the span's word count.

Substitutes are scored using asymmetric containment:

\[
\text{score}(A \to B) = \frac{|\text{contexts}(A) \cap \text{contexts}(B)|}{|\text{contexts}(A)|}
\]

This captures the fact that substitution is inherently asymmetric: ``dog'' can substitute for ``animal'' in many contexts, but not vice versa.

The CKY bottom-up process evaluates all binary decompositions of all spans, shortest first, caching substitution classes for reuse in longer spans. The minimum-energy decomposition at each level determines the parse structure.

\subsection{The Containment View}

The parser produces an interactive visualization called the \textit{Containment View}, which represents the parse hierarchy as nested circles:

\begin{itemize}
    \item Each multi-word span becomes a circle containing its children.
    \item Input words appear as nodes inside their parent span's circle.
    \item Substitute words appear alongside the input words they can replace.
    \item Sequence links (arrows) thread through the circles, showing how sequential data composes into hierarchical structure.
    \item Energy labels on each circle indicate the evidence strength of that decomposition level.
\end{itemize}

Crucially, substitutes from higher-level spans are decomposed through the parse tree: a 3-word substitute of ``found my hat'' is split according to the same binary structure, placing ``dropped'' near ``found'' and ``a pencil'' inside the ``my hat'' circle as ``a'' near ``my'' and ``pencil'' near ``hat''. Sequence links then connect these decomposed parts, threading through the nested circles and connecting to neighboring input words for context. Shared substitute words (appearing in multiple substitutes) are represented by a single node with multiple sequence links converging on it.

This visualization is central to the vision proposal: the nested circles of the Containment View map directly to the spatial containment of visual objects, and the sequence links map to the saccade path through the scene.

\section{From Words to Fixations: The Core Mapping}

The parser's architecture maps component-by-component to a visual perception system:

\begin{center}
\begin{tabular}{|l|l|}
\hline
\textbf{Language Parser} & \textbf{Visual Saccade Parser} \\
\hline
Word & Fixation patch (local visual sample) \\
Word sequence & Saccade scanpath \\
Text corpus & Eye-tracking corpus over many scenes \\
N-gram index & Fixation-subsequence index \\
Word substitution & Patch replacement in context \\
Context pattern (left/right words) & Surrounding fixation pattern \\
Containment score & Visual context overlap \\
Parse tree & Scene part-whole hierarchy \\
Containment circle & Visual object boundary \\
Sequence link & Saccade path through scene \\
Substitute node & Alternative visual element \\
\hline
\end{tabular}
\end{center}

The mapping is direct because the parser makes no assumptions specific to language. It operates on \textit{any} sequential data where (a)~elements can be compared for contextual similarity, (b)~a corpus of such sequences is available, and (c)~hierarchical structure is latent in the data. Visual fixation sequences satisfy all three conditions.

\section{Why Saccade Sequences Have Grammatical Structure}

\subsection{Scanpath Theory}

Noton and Stark's scanpath theory (1971) established that visual perception produces stereotypical fixation sequences. When viewing a particular visual pattern, observers execute characteristic saccade sequences that are repeated across viewings. These scanpaths are not random---they reflect the structure of the scene being viewed. Faces elicit characteristic scanpath patterns (eyes, nose, mouth in recurring sequences); reading a technical diagram follows the logical structure of the diagram.

This predictability is precisely what the parser exploits in language: the fact that word sequences are not random but follow structural patterns. The ``grammar'' of visual scenes is the set of regularities in fixation sequences, and energy minimization over substitutional evidence can discover it.

\subsection{Contrast with Top-Down Visual Grammars}

Zhu and Mumford's ``Stochastic Grammar of Images'' (2007) proposed a formal grammar for visual scenes, with production rules decomposing scenes into objects, objects into parts, and parts into primitives. Their And-Or graph representation captures both hierarchical decomposition (And-nodes) and alternative interpretations (Or-nodes).

However, their approach is fundamentally top-down: the grammar must be specified or learned from labelled examples. The production rules, non-terminal categories, and the structure of the And-Or graph are designed, not discovered. Our proposal inverts this: the ``grammar'' emerges bottom-up from substitutional evidence in raw fixation sequences. No categories, rules, or labelled training data are required. The And-nodes (hierarchical decomposition) emerge as the minimum-energy parse structure; the Or-nodes (alternatives) emerge as the substitution classes at each level.

\subsection{Distributional Semantics in Visual Scenes}

L\"{u}ddecke et al.\ (2019) demonstrated that distributional semantics---the principle that meaning is defined by context---extends from text to visual scenes. Objects that appear in similar spatial contexts in scenes develop similar distributional representations, just as words in similar linguistic contexts develop similar embeddings. Their work confirms that the substitutional principle underlying our parser has a direct visual analogue: visual elements that can replace each other in scene contexts are semantically related.

\section{Emergent Visual Objects as Substitution Classes}

In the parser, a ``cognitive object'' is not a fixed lexical entry but a substitution class---the set of words that can replace a given word in a specific structural context. The substitution class for ``hat'' in ``my hat'' differs from the substitution class for ``hat'' in ``wear a hat'' because different surrounding structures yield different substitutes.

In vision, this means: a ``visual object'' is not a fixed category (``chair'', ``car'', ``face'') but a substitution class of fixation patches. The set of patches that can replace a given patch in a specific fixation-sequence context \textit{defines} what that patch ``is'' in that context. A circular patch is a ``wheel'' when its surrounding fixation context includes car-body patches, a ``plate'' when the context includes table-and-food patches, and a ``clock face'' when the context includes wall-and-room patches.

This context-dependent object identity is a feature, not a limitation. It reflects how human perception actually works: the same retinal pattern is perceived differently depending on context.

\subsection{Asymmetric Containment in Vision}

The parser's containment scoring is asymmetric: $\text{score}(\text{dog} \to \text{animal})$ is high because most contexts where ``dog'' appears are also contexts where ``animal'' could appear, but $\text{score}(\text{animal} \to \text{dog})$ is low because ``animal'' appears in many contexts where ``dog'' would be inappropriate.

In vision, this asymmetry captures part-whole relationships naturally:

\begin{itemize}
    \item $\text{score}(\text{wheel} \to \text{car})$ is high: most fixation contexts containing a wheel patch also contain car patches.
    \item $\text{score}(\text{car} \to \text{wheel})$ is low: car contexts contain many patches beyond wheels.
    \item $\text{score}(\text{pupil} \to \text{eye})$ is high; $\text{score}(\text{eye} \to \text{pupil})$ is lower.
\end{itemize}

This asymmetry means that taxonomic and mereological (part-whole) relationships in visual scenes emerge from containment scoring without any explicit ontological engineering.

\subsection{Comparison with GLOM}

Hinton's GLOM architecture (2021) addresses the same fundamental problem: how to represent part-whole hierarchies for visual scenes in a neural network. GLOM uses ``islands of identical vectors'' at different levels of a columnar architecture to represent parse-tree nodes. Vectors at adjacent levels interact through bottom-up encoders and top-down decoders, with consensus (agreement between nearby columns) driving the emergence of part-whole structure.

There is a striking parallel: both GLOM and our proposal use agreement/consensus as the mechanism driving structural emergence. But they differ in a crucial respect:

\begin{itemize}
    \item \textbf{GLOM} uses a fixed number of levels in a designed architecture. The part-whole hierarchy must fit into the pre-specified number of levels.
    \item \textbf{Our proposal} discovers the number of hierarchical levels from the data. The CKY process builds structure bottom-up with as many levels as the evidence supports. A simple scene might have two levels; a complex scene might have five.
\end{itemize}

Moreover, GLOM operates on a single image. Our proposal operates on a \textit{corpus} of fixation sequences, drawing on the statistical regularity across many scenes to determine what constitutes a meaningful part-whole decomposition.

\section{The Containment View as Scene Structure}

The Containment View visualization produced by the parser represents hierarchical structure as nested circles with sequence links. For visual scene parsing, this representation takes on a direct spatial interpretation:

\textbf{Nested circles = spatial containment of visual objects.} A fixation-subsequence covering an eye region would form a circle containing pupil and iris patches. This circle would be nested inside a larger face circle, which would be nested inside a person circle. The nesting is not imposed---it emerges from the energy minimization over fixation-sequence evidence.

\textbf{Sequence links = saccade paths.} The arrows threading through the nested circles represent the actual saccade sequence: the path the eye took through the scene. Multiple substitute sequences show alternative saccade paths that would be contextually equivalent---other ways the eye could have traversed the same structural regions.

\textbf{Substitute nodes = visual alternatives.} Just as ``your'' and ``his'' appear as substitutes for ``my'' in the ``my hat'' circle, alternative visual patches would appear inside the corresponding object circle. A brown eye and a blue eye would both appear as substitutes inside the ``eye'' circle, connected by sequence links to the surrounding face context.

\textbf{Energy labels = evidence strength.} The energy on each containment circle indicates how well-supported that level of decomposition is by the fixation corpus. Low energy means many consistent substitutions exist at that level---the decomposition is well-attested. High energy means sparse evidence---the decomposition is uncertain. This provides a built-in confidence measure for each level of the visual hierarchy.

\section{Complementarity and Ambiguity}

Goertzel's framework proposes that evidence conservation governs inference, preventing the fabrication of support beyond premises. The parser embodies this: every structural claim must be grounded in corpus evidence. Extending this to vision means that visual structure cannot be ``hallucinated''---it must be supported by observed fixation statistics.

\subsection{Visual Ambiguity as Competing Parses}

Ambiguous visual stimuli have a natural interpretation in this framework. The Necker cube, which can be perceived as oriented in two different ways, would produce two parse structures with similar energies. Neither interpretation has decisively more substitutional support than the other, so both persist as competing minimum-energy decompositions. Perceptual switching between interpretations corresponds to alternation between parse structures.

The duck-rabbit figure would similarly produce two low-energy parses: one where the fixation subsequence over the ``beak/ears'' region has a substitution class drawn from bird contexts, and another where the same subsequence has a substitution class drawn from rabbit contexts. The ambiguity is precisely that the same fixation-sequence span supports two different substitution classes with comparable energy.

\subsection{Context-Dependent Complementarity}

In the parser, the word ``animals'' in ``cats are animals'' and ``animals eat grass'' develops different substitution classes---the first includes pets and predators, the second narrows to herbivores. The syllogistic inference ``cats eat grass'' fails because the shared term is not a single object but two context-dependent objects.

In vision, the analogous phenomenon is that the same visual pattern acquires different ``identities'' in different scene contexts. A vertical rectangle is a ``door'' in a building context and a ``book'' in a shelf context. The energy-minimizing process correctly treats these as different objects (different substitution classes) rather than forcing a single fixed categorization. This is Goertzel's complementarity made concrete: context-dependent measurement that prevents false transfer of properties between contexts.

\section{Advantages Over Existing Approaches}

\begin{enumerate}
    \item \textbf{No training labels.} Structure emerges from substitutional statistics over raw fixation sequences. No annotated datasets, bounding boxes, or category labels are required.

    \item \textbf{Context-dependent object identity.} Objects are not fixed categories but context-dependent substitution classes. The same visual pattern can be different ``objects'' in different scenes.

    \item \textbf{Naturally hierarchical.} The CKY-like bottom-up process discovers part-whole structure with as many levels as the evidence supports, without designing a fixed feature hierarchy or specifying the number of levels.

    \item \textbf{Anti-hallucination.} All structural claims are grounded in observed fixation evidence. The system cannot fabricate object structure not supported by the fixation corpus. This is Goertzel's evidence conservation applied to visual perception.

    \item \textbf{No pre-specified distance metric.} Visual similarity between patches is not computed in a fixed embedding space. It emerges from the energy landscape over decompositions---two patches are ``similar'' if they can substitute for each other in fixation-sequence context.

    \item \textbf{Cross-modal unification.} The same algorithm, the same energy function, and the same containment scoring apply to both language and vision. This suggests a domain-general mechanism for discovering hierarchical structure in sequential data, not a language-specific or vision-specific solution.
\end{enumerate}

\section{Implementation Sketch}

A concrete implementation would require the following components:

\textbf{Fixation encoding.} Each fixation in a scanpath is encoded as a feature vector describing the local visual patch. This could be a CNN patch descriptor, a histogram of oriented gradients, or even raw pixel statistics. The encoding must support similarity comparison.

\textbf{Discretization for indexing.} The parser's n-gram index requires discrete tokens. Fixation patch descriptors can be discretized by clustering (e.g., k-means on patch descriptors), assigning each fixation to its nearest cluster centroid. The cluster IDs become the ``vocabulary'' of visual tokens. This is analogous to visual word approaches in image retrieval, but applied to fixation sequences rather than spatial bag-of-words.

\textbf{Corpus construction.} A corpus of scanpath sequences is needed---many observers viewing many scenes. Existing eye-tracking datasets (MIT Saliency Benchmark, OSIE, SALICON) provide scanpath data over diverse scenes. The corpus would be indexed as fixation-token n-grams, exactly as the text parser indexes word n-grams.

\textbf{Containment scoring.} The asymmetric containment score transfers directly: for a candidate fixation-token $A$ and target token $B$, compute the fraction of $A$'s fixation-sequence contexts that are also contexts of $B$. This requires the same left-context and right-context pattern infrastructure as the text parser.

\textbf{CKY parsing.} Given a new scanpath (a sequence of fixation tokens), apply the same bottom-up energy minimization: enumerate all binary splits of all subsequences, compute substitution classes and energy for each, select minimum-energy decompositions, cache results for longer subsequences.

\textbf{Visualization.} The Containment View directly renders the discovered scene structure as nested circles with saccade-path links threading through.

\section{Discussion}

\subsection{Bridging the Semantic Gap}

A persistent challenge in computer vision is the ``semantic gap'' between low-level pixel representations and high-level semantic categories. The parser's multi-level substitution classes offer a principled bridge: at the lowest level, substitution classes group similar visual patches (texture categories); at intermediate levels, they group functional parts (object components); at the highest levels, they group whole objects and scene elements. Each level emerges from the energy minimization over the level below, with no explicit category engineering at any level.

\subsection{Cross-Modal Transfer}

If the same substitutional mechanism operates on both language and vision, there is a natural point of contact: the abstract substitution classes at higher levels of the hierarchy. A visual substitution class for ``things you sit on'' and a linguistic substitution class for ``chair, bench, stool, sofa'' may converge on the same structural role. This suggests a path toward grounded language understanding that does not require explicit symbol-grounding rules but emerges from shared substitutional structure across modalities.

\subsection{Active Perception}

The proposal is inherently compatible with theories of active perception and embodied cognition. Saccades are not passive recordings---they are active decisions about where to sample next. The energy landscape over decompositions could potentially guide saccade planning: the system could direct fixations toward regions where additional evidence would most reduce the energy of competing structural hypotheses, implementing an active perception strategy grounded in the same energy-minimization framework.

\subsection{Limitations}

Several challenges remain. The discretization of continuous fixation patches into discrete tokens loses information; the granularity of clustering affects what structure can be discovered. The computational cost of CKY parsing grows with sequence length, though the parser's caching and two-tier architecture mitigate this. The approach assumes that fixation-sequence regularities are rich enough to support substitutional analysis---an empirical question that can only be answered by implementation and testing.

\subsection{Conclusion}

We have proposed that the energy-minimizing substitutional mechanism underlying a working text parser can be applied to visual perception by treating saccade fixation sequences as the input data. The containment visualization produced by the parser maps directly to the spatial part-whole hierarchy of visual scenes. This proposal offers a unified, domain-general mechanism for discovering hierarchical structure in sequential data---grounded in evidence, with no pre-specified categories, distances, or grammar rules.

\vfill

\section*{References}

\begin{itemize}
    \item Noton, D. \& Stark, L. (1971). Scanpaths in eye movements during pattern perception. \textit{Science}, 171(3968), 308--311.
    \item Zhu, S.C. \& Mumford, D. (2007). A Stochastic Grammar of Images. \textit{Foundations and Trends in Computer Graphics and Vision}, 2(4), 259--362. \url{http://www.stat.ucla.edu/~sczhu/papers/Reprint_Grammar.pdf}
    \item Hinton, G. (2021). How to represent part-whole hierarchies in a neural network. \textit{arXiv:2102.12627}. \url{https://arxiv.org/abs/2102.12627}
    \item LeCun, Y. (2007). Energy-Based Models in Document Recognition and Computer Vision. \textit{ICDAR Keynote}. \url{https://yann.lecun.com/exdb/publis/pdf/lecun-icdar-keynote-07.pdf}
    \item Goertzel, B. (2026). Evidence is to Logic what Energy is to Dynamics. \url{https://bengoertzel.substack.com/p/evidence-is-to-logic-what-energy}
    \item L\"{u}ddecke, T., Agostini, A., Fauth, M., Tamosiunaite, M., \& W\"{o}rg\"{o}tter, F. (2019). Distributional semantics of objects in visual scenes in comparison to text. \textit{Artificial Intelligence}, 274, 44--65.
    \item Bruni, E., Tran, N.K. \& Baroni, M. (2014). Multimodal Distributional Semantics. \textit{Journal of Artificial Intelligence Research}, 49, 1--47.
    \item Freeman, R. (2026). Hamiltonian Dynamic Parser as a Foundation for Hypergraph Knowledge Representation. Technical analysis.
\end{itemize}

\end{document}
